\section{Analysis}
\label{sec:analysis}

In this section, we present a qualitative analysis to provide some insights how MAD works.
Unless otherwise stated, we report the overall results on the Common~MT dataset.

\begin{table}[t]
\centering
\begin{tabular}{l cc cc}
\toprule
\bf Method & \bf Bias$\downarrow$ & \bf Diversity$\uparrow$\\
\midrule
\textbf{Self-Reflect}  & 29.0 & 19.3\\
\textbf{MAD}                  & 24.8 & 49.7\\
\bottomrule
\end{tabular}
\vspace{-5pt}
\caption{Mitigation of Degeneration-of-Thought.}
\label{tab:mitigate-dot}
% \vspace{-5pt}
\end{table}



\subsection{Mitigation of DoT}

As mentioned in the Section~\ref{sec:introduction}, the DoT problem originates from three factors: (1) Bias and Distorted Perception, (2) Rigidity and Resistance to Change, and (3) Limited External Feedback. In our MAD framework, we introduce the views of other agents in the form of debates, solving the phenomenon of limited external feedback (problem 3). Next, this section will delve into the mitigation of problems 1 and 2 through experiments.

\begin{itemize}[leftmargin=10pt]
    \item \textbf{Bias}: We observe that LLMs often rely on direct intuition, which can lead to incorrect or inappropriate responses. To address this problem, we use human evaluation to determine the ambiguity error rate of LLMs' responses, examining if the LLM's output is biased.
    \item \textbf{Diversity}: LLMs are resistant to changing their answers and lack diverse reflection. The diversity of the translations is evaluated using the Self-BLEU score~\citep{yin2020meta}. In other words, methods lacking diverse reflection produce more similar translation candidates. Consequently, higher Self-BLEU scores mean lower diversity. We calculate text diversity via:
    \begin{equation}
    \scalebox{0.84}{$
    \mathrm{Diversity} = 100-\mathrm{Self\_BLEU}\left(Cand_1, Cand_2\right)
    $}
    \label{formula:diversity}
    \end{equation}
\end{itemize}

In formula~(\ref{formula:diversity}), candidates 1 and 2 represent the initial translation (base answer in Self-Reflection or affirmative side's response in MAD) and the current translation (possible modified answer after Self-Reflection or negative side's response in MAD).

As shown in Table~\ref{tab:mitigate-dot}, 
Bias and Rigidity are significant factors causing DoT. In addition, addressing these biases and stereotypes through self-reflection can be challenging. MAD framework effectively corrects inherent biases in translation, mitigates DoT, and considerably improves performance.

\subsection{Analysis of Judge}

In this section, we analyze the behavior of the judge for different settings of the debaters. 


\paragraph{Strong debaters with a weak judge work better than the reverse.}
To understand the roles of debaters and judge in MAD, we employ various combinations of models to initialize the agents. Specifically, we utilize the smaller language model (\texttt{vicuna-13b-v1.5-16k}) as a judge to evaluate the debate results of the more powerful LLMs (\texttt{GPT-3.5-Turbo}), and vice versa. 

The detailed experimental findings are presented in Table~\ref{tab:weak-judge}. The quality of the debaters' responses significantly impact the performance ceiling of MAD. Regardless of the model chosen for the judge, Turbo debaters consistently generate superior translations compared to Vicuna. In addition, the selection of the judge agent plays a secondary role. When Turbo debaters are involved, Vicuna, serving as the judge, underperforms Turbo across all test sets.

\begin{table}[t]
\centering
\begin{tabular}{l rr}
\toprule
\bf Judge LLM & \bf COMET & \bf HUMAN \\
\midrule
\multicolumn{3}{c}{\texttt{Vicuna-13b as Debaters}} \\
{\bf \texttt{Vicuna-13b}}  & 79.9 & 3.20 \\
{\bf \texttt{GPT-3.5-Turbo}}   & 80.4 & 3.25 \\
\midrule
\multicolumn{3}{c}{\texttt{GPT-3.5-Turbo as Debaters}} \\
{\bf \texttt{Vicuna-13b}}  & 83.2 & 3.47 \\
{\bf \texttt{GPT-3.5-Turbo}}   & 84.4 & 3.69 \\
\bottomrule
\end{tabular}
\vspace{-5pt}
\caption{Translation performance with different judge.}
\label{tab:weak-judge}
% \vspace{-12pt}
\end{table}


\paragraph{LLM may not act as an impartial judge when different LLMs are used as debaters.}
% \paragraph{LLM-based judge has a preference for outputs generated by itself.}
We study the behavior of agents by calculating how many times the judge chooses the answers of each debater as the final solution in different scenarios. The results are listed in Table~\ref{tab:behavior-agent} and we have the following observations:
\begin{itemize}[leftmargin=10pt]
    \item \textit{Same  LLM for All Agents} (Rows ~\textcircled{\small{1}} and \textcircled{\small{2}}): We find that the judge consistently favors the negative side, which is believed to contribute to the performance improvement in MAD. When encountering complex tasks, the affirmative side tends to make mistakes that should be corrected by the opposing side to achieve improvements.
    \item \textit{Debaters of Different LLMs} (Rows~\textcircled{\small{3}} and \textcircled{\small{4}}): We find that the judge shows a preference to the side with the same LLM as the backbone. This bias indicates that
LLMs might not be a fair judge~\citep{wang2023FairEval} when different LLMs are used for the agents.
\end{itemize}




\subsection{Analysis of Debaters}

In this section, we will discuss several factors of debaters that would affect the performance of MAD: \textit{debater number}, \textit{debate level}, and \textit{debate iteration}.


\paragraph{Increasing the number of debaters fails when backbone LLMs are poor at long-text modeling.}
It seems intuitive that increasing the number of debaters would enhance diversity of thought and subsequently improve performance. However, as shown in Table~\ref{tab:more-debaters}, an increase in the number of debaters has resulted in varying degrees of performance reduction.

\begin{figure}[t]
    \centering
    \includegraphics[height=0.67\linewidth]{imgs/Tit-for-tat.pdf}
    \vspace{-5pt}
    \caption{Translation performance with respect to the debate level on Lexical.}
    \label{fig:tit-for-tat}
    % \vspace{-5pt}
\end{figure}

\begin{table}[!t]
\centering
\setlength\tabcolsep{5pt}
\begin{tabular}{c ccc rrr}
\toprule
\multirow{2}{*}{\bf ID} & \multirow{2}{*}{\bf Jud} & \multicolumn{2}{c}{\bf Debater} & \multicolumn{3}{c}{\bf Winner}\\
\cmidrule(lr){3-4}\cmidrule(lr){5-7}
&   &  \bf Aff & \bf Neg &  \bf Aff & \bf Neg & \bf Tie  \\
\cmidrule(lr){1-4}\cmidrule(lr){5-7}
\textcircled{\small{1}} & \tt Turbo & \tt Turbo & \tt Turbo & 87 & \bf 104 & 9 \\
\textcircled{\small{2}} & \tt GPT-4 & \tt GPT-4 & \tt GPT-4 & 67 & \bf 124 & 9 \\
\hline
\textcircled{\small{3}} & \multirow{2}{*}{\tt GPT-4} & \tt Turbo & \tt GPT-4 & 52 & \bf 136 & 12 \\
\textcircled{\small{4}} & & \tt GPT-4 & \tt Turbo& \bf 120 & 77 & 3 \\
\bottomrule
\end{tabular}
\vspace{-5pt}
\caption{Number of times the judge chooses the answers of each debater based on different LLM.}
\label{tab:behavior-agent}
% \vspace{-10pt}
\end{table}

\begin{figure}
    \centering
    \includegraphics[height=0.67\linewidth]{imgs/distribution.pdf}
    \vspace{-5pt}
    \caption{Distribution of iteration rounds and a human score of each iteration subset.}
    \label{fig:distribution}
\end{figure}

\begin{table}
\centering
\begin{tabular}{l rr}
\toprule
\bf \# of Debaters & \small \bf COMET & \small \bf HUMAN \\
\midrule
\bf 2 (Default)   & 84.4 & 3.69 \\
\hdashline
\bf 3  & 83.1 & 3.58 \\
\bf 4  & 82.9 & 3.49 \\
\bottomrule
\end{tabular}
\vspace{-5pt}
\caption{Translation performance with more debaters.}
\label{tab:more-debaters}
% \vspace{-10pt}
\end{table}


To address this issue, we manually analyze the debate processes in approximately 10\% of the test subset. As the number of debaters increases, the length and complexity of the text also increase. Such LLM-based debaters tend to forget the views of other debaters during the debate. Moreover, it becomes more challenging for the judge to extract information from the debates for summarization. This suggests that the key challenge of MAD with more debaters lies in the limitations of the LLMs to handle long texts~\citep{liu2023lost}.



\paragraph{Appropriate "tit for tat" is beneficial for effective debate.}
We then study how the intensity of ``tit for tat'' affects the performance of MAD.
To achieve so, we design different instructions~(see Table~\ref{tab:tit-for-tat-prompt} in Appendix) to initialize the debaters' meta prompt. As shown in Figure~\ref{fig:tit-for-tat}, asking the debaters to ``tit for tat''~(i.e., higher disagreement) is necessary for MAD to achieve good performance. However, we find that ``\textit{must disagree with each other on every point }'' (with a disagreement of 0.988) does not lead to the best performance.
We speculate that continuous disagreement without finding common ground can contribute to polarization, where the debate becomes more about winning the argument than seeking truth or understanding. This can reinforce pre-existing biases and make it difficult to reach a meaningful consensus.


\paragraph{Complex questions require more iteration rounds of debate.}
In our experimental setup, we did not implement any additional stopping strategies besides setting the maximum debate iteration to 3. In other words, the judge can take an \textit{adaptive break} if it believes the optimal answer has already been obtained, efficiently ending the debate early.

To understand the distribution of iteration rounds and factors contributing to a longer debate process, we analyze the experimental results and present them in Figure~\ref{fig:distribution}. In the majority of cases, the optimal answer can be achieved through a single round of debate, demonstrating the efficiency of MAD. However, when translating more complex sentences (subsets with lower human scores), the judge requires additional iterations to gather adequate information from the debaters before making a final decision.
We also find that our MAD framework consistently brings performance improvements across all the three subsets, demonstrating its effectiveness.

\begin{figure}
    \centering
    \includegraphics[height=0.67\linewidth]{imgs/COMET-iteration.pdf}
    \vspace{-5pt}
    \caption{Performance with respect to the iteration of debate or self-reflection.}
    \label{fig:comet-iteration}
    % \vspace{-10pt}
\end{figure}


\paragraph{Adaptive break plays an important role to conclude the debate in the optimal moment.}
Intuitively, longer debates would encourage more diverse thinking. It raises the question of how the model's performance would be affected if constrained to conclude at a specific debate round.
For each iteration, we force the judge $J$ to extract the final answer ($a = J_e(H)$) instead of adaptively breaking the debate as in MAD. 

As shown in figure~\ref{fig:comet-iteration}, we can observe that MAD performs better than self-reflection as the iteration increases. However, the highest COMET score appears at the first iteration and is also lower than the result of the adaptive break. It indicates that, for most examples, MAD can generate good translations at the first iteration such that the debate should be stopped.
Forcing the debate to continue will harm the translation results, which demonstrates the reasonableness of our adaptive break strategy.    





